% Macro locali del problema; estratte da latex/SA_Book.tex.
\providecommand{\IniziaSudoku}[8]
{\psset{unit=3mm}
\pspicture(0,0)(8,8)
\psline(0,0)(0,8)(8,8)(8,0)(0,0)
\psline(0,4)(8,4)
\psline(4,0)(4,8)
\psline[linestyle=dotted,dotsep=1pt](2,0)(2,8)
\psline[linestyle=dotted,dotsep=1pt](6,0)(6,8)
\psline[linestyle=dotted,dotsep=1pt](0,2)(8,2)
\psline[linestyle=dotted,dotsep=1pt](0,6)(8,6)
\rput{0}(1,7){$#1$}
\rput{0}(3,7){$#2$}
\rput{0}(5,7){$#3$}
\rput{0}(7,7){$#4$}
\rput{0}(1,5){$#5$}
\rput{0}(3,5){$#6$}
\rput{0}(5,5){$#7$}
\rput{0}(7,5){$#8$}
}

\providecommand{\FinisciSudoku}[8]
{\rput{0}(1,3){$#1$}
\rput{0}(3,3){$#2$}
\rput{0}(5,3){$#3$}
\rput{0}(7,3){$#4$}
\rput{0}(1,1){$#5$}
\rput{0}(3,1){$#6$}
\rput{0}(5,1){$#7$}
\rput{0}(7,1){$#8$}
\endpspicture
}

\renewcommand{\Anno}{2005}

\renewcommand{\NumeroProblema}{1}

\renewcommand{\Testo}{\begin{minipage}[t]{12cm} Un piccolo Sudoku
è una tabella $4\times 4$, a sua volta suddivisa in quattro piccole tabelle $2\times 2$: ciascuna delle 16 caselle della tabella contiene una cifra (scelta tra 1, 2, 3, 4), in modo che nessuna delle 4 righe, delle 4 colonne, e delle 4 tabelle $2\times2$ contenga più di una volta la stessa cifra.  La tabella qui accanto
rappresenta un esempio di piccolo Sudoku.  \end{minipage}\hfill
	\begin{minipage}[t]{3cm}
		\mbox{}\vspace{-1ex}\\
		\hfill\IniziaSudoku{1}{2}{3}{4}{3}{4}{1}{2}
		\FinisciSudoku{2}{1}{4}{3}{4}{3}{2}{1}
	\end{minipage}
	
\begin{enumerate}
	\item Supponiamo che una delle diagonali del piccolo Sudoku contenga 4 cifre diverse.  Dimostrare che anche l'altra diagonale contiene 4 cifre diverse.
	
	\item Determinare quanti sono i piccoli Sudoku in cui anche le diagonali contengono 4 cifre diverse.
	
	\item  Determinare quanti sono i possibili piccoli Sudoku.
\end{enumerate}}

\renewcommand{\Risposta}{\begin{enumerate}
		\item[(a)]  \ldots
	
		\item[(b)]  48
	
		\item[(c)]  288
\end{enumerate}}

\renewcommand{\Aiutino}{\begin{enumerate}
	\renewcommand{\labelenumi}{(\alph{enumi})} 
	\item[(a)] Indichiamo con $A$, $B$, $C$, $D$ i 4 numeri sulla diagonale NW-SE. Quale numero può stare al secondo posto della prima riga?  Ci sono 2 possibilità.  In ciascuno dei 2 casi, il piccolo Sudoku si completa in maniera unica.  In entrambi i casi tutti i numeri appaiono anche sull'altra diagonale.

	\item[(b)]  Abbiamo appena visto che ci sono 2 casi, oltre alla possibilità di permutare i simboli.

	\item[(c)]  Indichiamo con $A$, $B$, $C$, $D$ i 4 numeri nel primo quadrato $2\times 2$. Scambiando eventualmente la terza e la quarta riga, e scambiando eventualmente la terza e la quarta colonna, possiamo assumere che nella terza casella della diagonale NW-SE ci sia una $A$. A questo punto nell'ultimo posto libero della diagonale ci possono essere $B$, $C$ o $D$, ed in ogni caso il piccolo Sudoku si completa in modo unico.
\end{enumerate}}

\renewcommand{\SoluzioneUfficiale}{\TitoloSoluzione \DomandaA Per simmetria 
possiamo supporre che 4 numeri distinti (che indicheremo con $A$, $B$, $C$, $D$) stiano sulla diagonale NW-SE.
\begin{center}
	\IniziaSudoku{A}{}{}{}{}{B}{}{}
	\FinisciSudoku{}{}{C}{}{}{}{}{D}
\end{center}

Nel primo quadrato, a destra della $A$, può esserci una $C$ o una $D$. In entrambi i casi, il Sudoku si completa facilmente in modo unico, come indicato qui sotto (si noti che, come avremmo potuto aspettarci, i due Sudoku ottenuti sono l'uno il simmetrico dell'altro rispetto alla diagonale considerata).
\begin{center}
	\hfill
	\IniziaSudoku{A}{C}{D}{B}{D}{B}{A}{C}
	\FinisciSudoku{B}{D}{C}{A}{C}{A}{B}{D}
	\hfill
	\IniziaSudoku{A}{D}{B}{C}{C}{B}{D}{A}
	\FinisciSudoku{D}{A}{C}{B}{B}{C}{A}{D}
	\hfill\mbox{}
\end{center}

Poiché in entrambi i casi anche la diagonale NE-SW contiene tutti i simboli, la tesi è dimostrata.

\DomandaB Ci sono $4!=24$ modi di disporre 4 numeri diversi sulla
diagonale NW-SE. Fissati questi, dal punto precedente sappiamo che ci sono 2 modi di completare il Sudoku.  Pertanto i possibili Sudoku con 4 numeri diversi sulle diagonali sono $24\cdot 2=48$.

\DomandaC Indichiamo con $A$, $B$, $C$, $D$ i numeri del primo
quadrato $2\times 2$, come in figura.
\begin{center}
	\IniziaSudoku{A}{B}{}{}{C}{D}{}{}
	\FinisciSudoku{}{}{}{}{}{}{}{}
\end{center}

Ci sono ora 4 possibili posizioni per la lettera $A$ all'interno del quadrato $2\times 2$ in basso a destra. Tuttavia, per non dover esaminare tutti e 4 i casi, possiamo eventualmente scambiare tra di loro la terza e quarta colonna, ed eventualmente scambiare tra di loro la terza e quarta riga, in modo che la $A$ venga a trovarsi in alto a sinistra in tale quadrato. Si noti che gli scambi che abbiamo eventualmente fatto mantengono la proprietà dei piccoli Sudoku, per cui se partiamo da un piccolo Sudoku alla fine ci troviamo ancora un piccolo Sudoku. Completando nei 3 modi possibili la diagonale, ci siamo quindi ricondotti ad esaminare i seguenti 3 casi.
\begin{center}
	\hfill
	\IniziaSudoku{A}{B}{}{}{C}{D}{}{}
	\FinisciSudoku{}{}{A}{}{}{}{}{B}
	\hfill
	\hfill
	\IniziaSudoku{A}{B}{}{}{C}{D}{}{}
	\FinisciSudoku{}{}{A}{}{}{}{}{C}
	\hfill
	\IniziaSudoku{A}{B}{}{}{C}{D}{}{}
	\FinisciSudoku{}{}{A}{}{}{}{}{D}
	\hfill\mbox{}
\end{center}

È ora facile vedere che tutti quanti si possono completare in modo unico come segue (conviene partire piazzando ovunque possibile le lettere che abbiamo già messo nel quadrato $2\times 2$ in fondo a destra).
\begin{center}
	\hfill
	\IniziaSudoku{A}{B}{D}{C}{C}{D}{B}{A}
	\FinisciSudoku{B}{C}{A}{D}{D}{A}{C}{B}
	\hfill
	\IniziaSudoku{A}{B}{C}{D}{C}{D}{B}{A}
	\FinisciSudoku{D}{C}{A}{B}{B}{A}{D}{C}
	\hfill
	\IniziaSudoku{A}{B}{D}{C}{C}{D}{B}{A}
	\FinisciSudoku{D}{C}{A}{B}{B}{A}{C}{D}
	\hfill\mbox{}
\end{center}

In conclusione, abbiamo 4 possibili posizioni per la $A$ all'interno del quadrato $2\times 2$ in fondo a destra, 3 modi di concludere una volta fissata tale $A$, più ovviamente le 24 possibili permutazioni dei simboli: i possibili piccoli Sudoku sono pertanto
$3\cdot 4\cdot 24=288$.}

\renewcommand{\SoluzioniAlternative}{\TitoloSoluzioneAlternativaA Per 
il punto (c), esistono ovviamente molti altri modi di contare i piccoli Sudoku. Una possibilità è la seguente. Siano $A$, $B$, $C$, $D$ i numeri della prima riga. Abbiamo allora 4 possibili modi di completare la seconda riga.
\begin{center}
	\hfill
	\IniziaSudoku{A}{B}{C}{D}{C}{D}{A}{B}
	\FinisciSudoku{}{}{}{}{}{}{}{}
	\hfill
	\IniziaSudoku{A}{B}{C}{D}{D}{C}{B}{A}
	\FinisciSudoku{}{}{}{}{}{}{}{}
	\hfill
	\IniziaSudoku{A}{B}{C}{D}{C}{D}{B}{A}
	\FinisciSudoku{}{}{}{}{}{}{}{}
	\hfill
	\IniziaSudoku{A}{B}{C}{D}{D}{C}{A}{B}
	\FinisciSudoku{}{}{}{}{}{}{}{}
	\hfill\mbox{}	
\end{center}

Consideriamo il primo, ed esaminiamo il quadrato $2\times 2$ in basso a sinistra. La prima colonna deve contenere $B$ e $D$, la seconda $A$ e $C$: abbiamo quindi 4 possibilità, che si completano tutte in modo unico come segue.
\begin{center}
	\hfill
	\IniziaSudoku{A}{B}{C}{D}{C}{D}{A}{B}
	\FinisciSudoku{B}{A}{D}{C}{D}{C}{B}{A}
	\hfill
	\IniziaSudoku{A}{B}{C}{D}{C}{D}{A}{B}
	\FinisciSudoku{B}{C}{D}{A}{D}{A}{B}{C}
	\hfill
	\IniziaSudoku{A}{B}{C}{D}{C}{D}{A}{B}
	\FinisciSudoku{D}{A}{B}{C}{B}{C}{D}{A}
	\hfill
	\IniziaSudoku{A}{B}{C}{D}{C}{D}{A}{B}
	\FinisciSudoku{D}{C}{B}{A}{B}{A}{D}{C}
	\hfill\mbox{}	
\end{center}

Ragionando allo stesso modo a partire dalla seconda possibilità, otteniamo i seguenti 4 piccoli Sudoku.
\begin{center}
	\hfill
	\IniziaSudoku{A}{B}{C}{D}{D}{C}{B}{A}
	\FinisciSudoku{B}{A}{D}{C}{C}{D}{A}{B}
	\hfill
	\IniziaSudoku{A}{B}{C}{D}{D}{C}{B}{A}
	\FinisciSudoku{B}{D}{A}{C}{C}{A}{D}{B}
	\hfill
	\IniziaSudoku{A}{B}{C}{D}{D}{C}{B}{A}
	\FinisciSudoku{C}{D}{A}{B}{B}{A}{D}{C}
	\hfill
	\IniziaSudoku{A}{B}{C}{D}{D}{C}{B}{A}
	\FinisciSudoku{C}{A}{D}{B}{B}{D}{A}{C}
	\hfill\mbox{}	
\end{center}

Consideriamo ora la terza possibilità. Ancora una volta abbiamo 4 possibili modi di riempire il quadrato $2\times 2$ in basso a sinistra, ma di questi 2 si completano in modo unico e 2 non si possono completare.
\begin{center}
	\hfill
	\IniziaSudoku{A}{B}{C}{D}{C}{D}{B}{A}
	\FinisciSudoku{B}{A}{D}{C}{D}{C}{A}{B}
	\hfill
	\IniziaSudoku{A}{B}{C}{D}{C}{D}{B}{A}
	\FinisciSudoku{D}{C}{A}{B}{B}{A}{D}{C}
	\hfill
	\IniziaSudoku{A}{B}{C}{D}{C}{D}{B}{A}
	\FinisciSudoku{B}{C}{?}{?}{D}{A}{?}{?}
	\hfill
	\IniziaSudoku{A}{B}{C}{D}{C}{D}{B}{A}
	\FinisciSudoku{D}{A}{?}{?}{B}{C}{?}{?}
	\hfill\mbox{}	
\end{center}

Discorso analogo per la quarta possibilità.

\begin{center}
	\hfill
	\IniziaSudoku{A}{B}{C}{D}{D}{C}{A}{B}
	\FinisciSudoku{B}{A}{D}{C}{C}{D}{B}{A}
	\hfill
	\IniziaSudoku{A}{B}{C}{D}{D}{C}{A}{B}
	\FinisciSudoku{C}{D}{B}{A}{B}{A}{D}{C}
	\hfill
	\IniziaSudoku{A}{B}{C}{D}{D}{C}{A}{B}
	\FinisciSudoku{B}{D}{?}{?}{C}{A}{?}{?}
	\hfill
	\IniziaSudoku{A}{B}{C}{D}{D}{C}{A}{B}
	\FinisciSudoku{C}{A}{?}{?}{B}{D}{?}{?}
	\hfill\mbox{}	
\end{center}

Abbiamo quindi ottenuto 12 piccoli Sudoku, da moltiplicare per le 24 permutazioni dei simboli, per un totale di $12\cdot 24=288$ piccoli
Sudoku.}

\renewcommand{\Osservazioni}{
}

\renewcommand{\Commenti}{
}